\documentclass[11pt]{article}

\usepackage[utf8]{inputenc}
\usepackage[T1]{fontenc}
\usepackage{amsmath}
\usepackage{amssymb}
\usepackage{booktabs}
\usepackage{hyperref}
\usepackage{microtype}
\usepackage{geometry}

\geometry{margin=1in}

\title{Zero-Training Acquisition and Multi-Agent Transmission of a Novel Epistemic Protocol Across Heterogeneous Language Models}

\author{
  Iva Vrtari\'c \\
  Independent Researcher
}

\date{}

\begin{document}

\maketitle

\begin{abstract}
ΛLingua is a compact symbolic protocol for representing assertions, evidence, uncertainty, hypotheses, and support relations across heterogeneous language models. We study whether unmodified models can acquire the protocol from a frozen specification and small in-context curriculum, how much teaching is required before competence becomes stable, and whether epistemic structure survives model-to-model transmission. Experiment Zero establishes a zero-training baseline across multiple local model families and shows that protocol-looking output can diverge sharply from semantic correctness. Experiment One introduces disjoint checkpoint exams, corrective teaching rounds, cross-domain transfer, and bootstrap-compression conditions to measure acquisition cost and epistemic-status fidelity. A later Agent-Wire relay study will test semantic drift across heterogeneous multi-model handoffs. This manuscript is a work in progress; acquisition and relay results will be inserted only after the corresponding frozen experiments are complete.
\end{abstract}

\section{Introduction}

Language models increasingly exchange summaries, tool results, intermediate states, and delegated tasks. In these handoffs, lexical content may survive while epistemic status changes. A possibility may become an assertion, evidence may detach from the claim it originally supported, or uncertainty may disappear during compression.

ΛLingua is designed as a deliberately small protocol for making those distinctions explicit.

The frozen ΛLingua v0 epistemic core is:

\begin{center}
\begin{tabular}{ll}
\toprule
Symbol & Meaning \\
\midrule
$\phi$ & assertion \\
$\epsilon$ & evidence \\
? & unknown \\
$\hat{h}$ & hypothesis \\
$\because$ & supported-by relation \\
\bottomrule
\end{tabular}
\end{center}

A canonical example is:

\begin{verbatim}
φ:HTTP503 ∵ ε:E1
\end{verbatim}

The marker $\phi$ denotes the represented epistemic role of the proposition; it does not certify objective truth.

\section{Research Questions}

The project is organized around three questions:

\begin{enumerate}
  \item Can heterogeneous, unmodified language models acquire ΛLingua from a compact frozen specification?
  \item How much in-context teaching is required before acquisition becomes stable on unseen material?
  \item Can a chain of heterogeneous models preserve shared epistemic state across repeated handoffs?
\end{enumerate}

\section{Experiment Zero: Zero-Training Baseline}

Experiment Zero tested zero-training pickup under frozen bootstrap conditions.

The initial model families were Gemma 3 1B, Qwen 3 1.7B, and Qwen 3.5 2B. A later replication added Phi-4 Mini 3.8B.

The experiment distinguishes surface protocol resemblance from semantic preservation. Outputs are evaluated for semantic fidelity, provenance survival, uncertainty survival, and protocol violations.

\subsection{Illustrative failure}

For one task, the canonical state was:

\begin{verbatim}
φ:HEALTH_CHECK_FAILED ∵ ε:E7
ε:E8
?:DNS_CAUSED_INCIDENT
\end{verbatim}

A Phi replication produced:

\begin{verbatim}
φ:DNS_LOOKUP ∵ ε:E8
?:HEALTH_CHECK_FAILURE
ĥ:DNS_CAUSE
\end{verbatim}

The output is syntactically recognizable as ΛLingua while reassigning epistemic roles and provenance. This distinction motivates explicit epistemic-status scoring.

\section{Experiment One: Acquisition}

Experiment One moves from zero-training pickup to explicit in-context acquisition.

The curriculum is evaluated at five checkpoints:

\begin{verbatim}
C0   specification only
C1   + worked demonstrations
C2   + corrective round 1
C3   + corrective round 2
C4   + corrective round 3
\end{verbatim}

Each checkpoint uses a separate structurally matched exam form:

\begin{verbatim}
C0 -> X0
C1 -> X1
C2 -> X2
C3 -> X3
C4 -> X4
\end{verbatim}

The forms are content-disjoint while sharing matched epistemic skeletons. This prevents repeated exposure to one exam from becoming part of the teaching process.

\subsection{Acquisition metrics}

In addition to the Experiment Zero metrics, Experiment One introduces Epistemic Status Fidelity (ESF) and acquisition-token measures.

\begin{description}
  \item[AT$\Lambda_{first}$] the first checkpoint at which the frozen competence threshold is reached.
  \item[AT$\Lambda_{stable}$] the first checkpoint from which competence remains satisfied through C4.
\end{description}

ATΛ counts teaching tokens only.

\subsection{Cross-domain transfer}

Transfer items preserve the epistemic structures while changing subject matter across infrastructure, science, and logistics domains.

\subsection{Bootstrap compression}

The frozen compression study includes full specification, nested deletion variants, examples-only, and symbols-only conditions.

\section{Multi-Agent Transmission}

The final planned study places ΛLingua inside a controlled Agent-Wire relay.

A semantic state will be passed through heterogeneous models while every handoff is preserved as a durable trace. The study will measure semantic drift, provenance decay, uncertainty laundering, hypothesis inflation, evidence mutation, error amplification, and repair behavior.

\section{Reproducibility}

Experimental inputs are frozen before inference. The repository preserves versioned specifications, benchmark files, SHA-256 manifests, disjoint exam forms, leakage checks, pinned model identities, fresh evaluation contexts, raw model responses, and deterministic scoring definitions.

Changes to a frozen protocol, benchmark, threshold, or inference condition create a new experiment version rather than rewriting the original experiment.

\section{Results}

\textbf{Work in progress.}

Experiment Zero and the Phi replication are complete. Experiment One official acquisition inference and the later Agent-Wire relay study must be completed before this section is finalized.

\section{Discussion}

\textbf{Placeholder.}

This section will interpret model-family differences only after the full acquisition results are available.

\section{Limitations}

The current experiments use small local model families and a deliberately compact protocol. Results should not be generalized to larger proprietary systems without direct testing. The study evaluates preservation of represented epistemic state, not objective truth.

\section{Conclusion}

ΛLingua treats epistemic transport as a measurable problem. The project asks how little symbolic structure is required for heterogeneous models to preserve claims, evidence, uncertainty, and hypotheses across acquisition and communication.

\bibliographystyle{plain}
\bibliography{references}

\end{document}
